% Part: sets-functions-relations
% Chapter: arithmetization
% Section: reflections

\documentclass[../../../include/open-logic-section]{subfiles}

\begin{document}

\olfileid[ar]{sfr}{arith}{ref}
\olsection{بعض التأملات الفلسفية}

هذا آخر ما قصدناه من المسائل التقنية؛ فلننظر الآن فيما أفدناه منها.

أفدنا منها، فيما لا يكاد ينازع فيه، رياضيات بحتة {جميلة}. وفوق ذلك حصل تحقيق عميق لبعض المفاهيم: فإن إدراك أن خاصية الاكتمال هي الفارق الحاسم بين الحقيقية والنسبية نفاذ بعيد في النظر. ثم إن بناء الحقيقية صراحة بوصفها قطوع ديديكند يثبت أساس موضوع التحليل؛ فمفهوم \emph{الحقل المرتب الكامل} متسق المعنى، لأن القطوع تقدم لنا حقلًا بهذه الصفات نفسها.

على أن هذه المكاسب لا تمنع من إبداء تحفظات في دلالتها.

أولها أن النظر إلى الحقيقية على أنها قطوع لا يظهر أنه \emph{يزيد} صرامة على النظر إليها من خلال تمثيلاتها العشرية المألوفة، وإن امتدت إلى غير نهاية. ولهذا «البناء» العشري شبه ببناء الحقيقية بواسطة متتاليات كوشي؛ على أن أصله كان معروفًا عند الرياضيين منذ أوائل القرن السابع عشر (انظر \olref[cauchy]{sec}). فالزيادة الحقيقية في الصرامة ناشئة عن إدراك خاصية الاكتمال؛ أما القدرة على تمثيل الأعداد الحقيقية بمجموعات معينة، فلعلها لا تستحق في ذاتها ذلك القدر من الاهتمام.

وأخفى من ذلك وجه زيادة الصرامة بالبناء الحسابي للصحيحة أو النسبية، مع أن بناءهما أيسر كثيرًا. وهنا ملاحظة جديرة بالتسجيل: ما إن \emph{نبني} الصحيحة فئات تكافؤ لأزواج طبيعية، والنسبية فئات تكافؤ لأزواج صحيحة، والحقيقية مجموعات من النسبية، حتى \emph{نترك هذه الأبنية وراءنا}. فلا يرغب أحد في \emph{استدعائها} أثناء البرهان الرياضي، إلا إذا كان البرهان نفسه لتحقيق أنها تعمل كما ينبغي. فالحديث مباشرة عن عدد حقيقي أيسر كثيرًا من الحديث عن مجموعة من مجموعات من مجموعات من مجموعات من مجموعات من مجموعات من مجموعات من الأعداد الطبيعية.
%(For reals were constructed as sets of rationals; and these were constructed as equivalence classes of ordered pairs of integers, i.e., sets of sets of sets of integers; etc.).
%2/3 = [2,3]
%	i.e. {<2,3>, ... }
%	{ {{2},{2,3}}, ... }
%
%reals are 
%	sets of rationals 
%
%rationals are 
%	sets of sets of sets of integers
%
%integers are
%	sets of sets of sets of naturals

وأشد هذه الدعاوى تعرضًا للشك أن تكشف التعريفات عما \emph{تكونه} الصحيحة والنسبية والحقيقية \emph{في نفسها من جهة الميتافيزيقا}. فليس بيّنًا أن الحقيقية، مثلًا، \emph{هي} تلك المجموعات المعينة من مجموعات من مجموعات\ldots. وأصل الإشكال أن البناء يمكن إجراؤه على وجوه كثيرة مختلفة. ففي الحقيقية أبنية ذات اختلاف حقيقي مهم (انظر \olref[cauchy]{sec})؛ بل يمكن تحصيل اختلاف في البناء بأمر أيسر من ذلك جدًّا: نغير تعريف الزوج المرتب قليلًا، كما في \olref[sfr][rel][ref]{sec}. ولو أخذنا بذلك التعريف البديل لعملت أبنيتنا كما تعمل الآن، ولكن الكيانات الناتجة لكانت غير هذه. فثمة أبنية متنافسة كثيرة للصحيحة والنسبية والحقيقية في نظرية المجموعات. ومع هذا التعدد، يكون الجزم بأن الأعداد الصحيحة مثلًا \emph{هذه} المجموعات دون \emph{تلك} اختيارًا اعتباطيًّا محرجًا. وهذا، كما في \olref[sfr][rel][ref]{sec}، مثال للحجة التي اشتهرت بـ\citealt{Benacerraf1965}.

وثمة وجه آخر للنظر، فإن في أبنيتنا أمرًا \emph{عجيبًا}. بدأنا بالطبيعية، ثم بنينا الصحيحة، وبنينا «${\OLMathNumeral{0}}$ الصحيح»، وهو $ \equivrep{{\OLMathNumeral{0}},{\OLMathNumeral{0}}}{\Intequiv}$. لكن ${\OLMathNumeral{0}} \neq
\equivrep{{\OLMathNumeral{0}},{\OLMathNumeral{0}}}{\Intequiv}$؛ بل ليس بحسب هذا البناء \emph{أي} عدد طبيعي عددًا صحيحًا. وهذا بعيد جدًّا عن التصور المعتاد. وقد ذكرنا في \olref[sfr][set][imp]{sec}، من غير إطالة في الاحتجاج، أن $\Nat \subseteq \Rat$. فلو كانت الأبنية تحدد بالضبط \emph{ماهية} الأعداد، لكان ذلك القول ظاهر البطلان.

فإذا جمعنا أطراف النظر، فما النتيجة؟ إننا، بالعمل في نظرية مجموعات ساذجة وافتراض الطبيعية معطاة لنا، نقدر على \emph{معاملة} الصحيحة والنسبية والحقيقية معاملة مجموعات معينة. ومن هذه الجهة نستطيع \emph{تضمين} نظريات تلك الكيانات في نظرية مجموعات. غير أن ما يترتب على هذا التضمين فلسفيًّا ليس بيّنًا بهذه السهولة.

وليس هذا كله فصل الخطاب!{} وإنما المراد أن القول بأن البناء الحسابي للأعداد الحقيقية \emph{له} دلالة فلسفية عميقة يحتاج، فوق البناء الرياضي، إلى حجة \emph{فلسفية} أخرى.

%Additionally: for the entire duration of this chapter, we have treated
%the natural numbers as simply \emph{given} to us. But what can we do
%with them? That is the subject matter for the next chapter.

\end{document}
